\documentclass[11pt,oneside]{article}

\pdfminorversion=7
\usepackage{QuizStyle}

\hypersetup{
    pdftitle={20261001 Quiz for MATH261 Discrete Mathematics},
    pdfauthor={Jungwoo Kim},
    pdfsubject={Quiz for MATH261 Discrete Mathematics},
    pdfcreator={LaTeX}
}

\begin{document}

\quiztitle{20261001 Quiz}

\begin{problem}{1}
Using mathematical induction, prove that
\[
    2^n\geq n^2
\]
for every $n\in\N$ with $n\geq4$.
\end{problem}

\begin{solution}
For $n=4$,
\[
    2^4=16=4^2,
\]
so the statement holds in the base case.

Assume that, for some $m\geq4$,
\[
    2^m\geq m^2.
\]
Then
\[
    2^{m+1}=2\cdot2^m\geq2m^2.
\]
Since $m\geq4$,
\[
    2m^2-(m+1)^2
    =m^2-2m-1
    =(m-1)^2-2>0.
\]
Hence
\[
    2^{m+1}\geq2m^2>(m+1)^2.
\]
Therefore, by mathematical induction,
\[
    2^n\geq n^2
\]
for every $n\in\N$ with $n\geq4$.
\end{solution}

\begin{problem}{2}
Let $n,k\in\Nzero$ with $n\geq k$. Prove that
\[
    \sum_{i=k}^{n}\binom{i}{k}
    =\binom{n+1}{k+1}.
\]
\end{problem}

\begin{solution}[1]
Count the $(k+1)$-element subsets of $[n+1]$ according to their largest element.
If the largest element is $i+1$, where $k\leq i\leq n$, then the other $k$ elements must be chosen from $[i]$, giving
\[
    \binom{i}{k}
\]
choices. These cases are disjoint and exhaustive. Therefore,
\[
    \sum_{i=k}^{n}\binom{i}{k}
    =\binom{n+1}{k+1}.
\]
\end{solution}

\begin{solution}[2]
Fix $k\in\Nzero$ and use induction on $n\geq k$.
For $n=k$,
\[
    \sum_{i=k}^{k}\binom{i}{k}
    =\binom{k}{k}
    =1
    =\binom{k+1}{k+1}.
\]

Assume that, for some $m\geq k$,
\[
    \sum_{i=k}^{m}\binom{i}{k}
    =\binom{m+1}{k+1}.
\]
Then, by Pascal's identity,
\[
\begin{aligned}
    \sum_{i=k}^{m+1}\binom{i}{k}
    &=\sum_{i=k}^{m}\binom{i}{k}+\binom{m+1}{k}\\
    &=\binom{m+1}{k+1}+\binom{m+1}{k}\\
    &=\binom{m+2}{k+1}.
\end{aligned}
\]
Hence the identity holds for every $n\geq k$.
\end{solution}

\begin{remark}
This identity is commonly called the \emph{hockey-stick identity}, since the corresponding entries in Pascal's triangle form a hockey-stick pattern.
\end{remark}

\begin{problem}{3}
Show that every postage amount of at least $12$ cents can be formed using only $3$-cent and $7$-cent stamps.
\end{problem}

\begin{solution}[1]
Let $n\geq12$. We distinguish the three possible residues of $n$ modulo $3$.
\begin{problemchoices}
    \item If $n\equiv0\pmod3$, then $n$ can be formed using only $3$-cent stamps.
    \item If $n\equiv1\pmod3$, then $n\geq13$ and
    \[
        n=7+(n-7),
    \]
    where $n-7$ is a nonnegative multiple of $3$.
    \item If $n\equiv2\pmod3$, then $n\geq14$ and
    \[
        n=14+(n-14),
    \]
    where $n-14$ is a nonnegative multiple of $3$.
\end{problemchoices}
Thus every $n\geq12$ can be formed from $3$-cent and $7$-cent stamps.
\end{solution}

\begin{solution}[2]
Use generalized induction. The first three cases are
\[
    12=3+3+3+3,
    \qquad
    13=3+3+7,
    \qquad
    14=7+7.
\]
Now let $m\geq14$ and assume that the postage amounts $m-2$, $m-1$, and $m$ can all be formed. In particular, $m-2$ can be formed, so adding one $3$-cent stamp gives
\[
    m+1=(m-2)+3.
\]
Hence $m+1$ can also be formed. Therefore every postage amount of at least $12$ cents can be formed.
\end{solution}

\begin{remark}
This is a special case of the \emph{Frobenius coin problem}. For relatively prime positive integers $a$ and $b$, the largest nonnegative integer not representable as
\[
    ax+by,
    \qquad x,y\in\Nzero,
\]
is the Frobenius number $g(a,b)$. In the two-denomination case, Sylvester's formula gives
\[
    g(a,b)=ab-a-b.
\]
Thus $g(3,7)=11$, which explains the threshold $12$ in this problem.
\end{remark}

\begin{problem}{4}
Let $d\in\N$. Define a relation $R$ on $\Z$ by
\[
    xRy
    \quad\Longleftrightarrow\quad
    d\mid(x-y).
\]
Show that $R$ is an equivalence relation and determine its equivalence classes.
\end{problem}

\begin{solution}
For every $x\in\Z$,
\[
    d\mid(x-x)=0,
\]
so $R$ is reflexive.

If $xRy$, then $d\mid(x-y)$. Since $y-x=-(x-y)$, we have $d\mid(y-x)$, so $yRx$. Thus $R$ is symmetric.

If $xRy$ and $yRz$, then $d\mid(x-y)$ and $d\mid(y-z)$. Hence
\[
    d\mid\bigl((x-y)+(y-z)\bigr)=x-z,
\]
so $xRz$. Thus $R$ is transitive, and therefore $R$ is an equivalence relation.

For $a\in\Z$,
\[
\begin{aligned}
    [a]_R
    &=\{x\in\Z\mid d\mid(x-a)\}\\
    &=\{a+qd\mid q\in\Z\}.
\end{aligned}
\]
Thus the distinct equivalence classes are
\[
    [0]_R,[1]_R,\ldots,[d-1]_R.
\]
These are precisely the residue classes modulo $d$, and $xRy$ is the congruence relation $x\equiv y\pmod d$.
\end{solution}

\begin{problem}{5}
Prove the following identities. We use the convention $\binom{a}{b}=0$ when $a,b\in\Nzero$ and $b>a$.
\begin{problemsubparts}
    \item For $n,m\in\Nzero$ with $n\geq m$,
    \[
        \sum_{k=0}^{m}\binom{n-k}{m-k}
        =\binom{n+1}{m}.
    \]

    \item For $u,v,m\in\Nzero$,
    \[
        \sum_{k=0}^{m}\binom{u}{k}\binom{v}{m-k}
        =\binom{u+v}{m}.
    \]

    \item For $n,m\in\Nzero$ with $n\geq m$,
    \[
        \sum_{k=m}^{n}\binom{k}{m}\binom{n}{k}
        =\binom{n}{m}2^{n-m}.
    \]

    \item For $n\in\N$ and $0\leq m\leq n$,
    \[
        \sum_{k=0}^{m}(-1)^k\binom{n}{k}
        =(-1)^m\binom{n-1}{m}.
    \]
\end{problemsubparts}
\end{problem}

\begin{solution}[1]
\begin{problemsubparts}
    \item Count the $m$-element subsets of $[n+1]$ according to the first element of $[m+1]$ that is omitted. If $k+1$ is the first omitted element, then $1,\ldots,k$ are selected, $k+1$ is not selected, and the remaining $m-k$ elements are chosen from
    \[
        \{k+2,\ldots,n+1\},
    \]
    which has $n-k$ elements. Hence this case contributes
    \[
        \binom{n-k}{m-k}.
    \]
    Since an $m$-element subset cannot contain all of $[m+1]$, these cases are disjoint and exhaustive. Therefore
    \[
        \sum_{k=0}^{m}\binom{n-k}{m-k}
        =\binom{n+1}{m}.
    \]

    \item Let $U$ and $V$ be disjoint sets with $|U|=u$ and $|V|=v$. Count the $m$-element subsets of $U\cup V$ according to the number $k$ of selected elements from $U$. For a fixed $k$, there are
    \[
        \binom{u}{k}\binom{v}{m-k}
    \]
    choices. Summing over $k$ gives the left-hand side, while choosing the $m$ elements directly from $U\cup V$ gives the right-hand side. Hence
    \[
        \sum_{k=0}^{m}\binom{u}{k}\binom{v}{m-k}
        =\binom{u+v}{m}.
    \]

    \item Count the pairs $(M,K)$ satisfying
    \[
        M\subseteq K\subseteq[n],
        \qquad
        |M|=m.
    \]
    If $|K|=k$, first choose $K$ and then choose $M\subseteq K$, giving
    \[
        \binom{n}{k}\binom{k}{m}
    \]
    choices. Summing over $k=m,\ldots,n$ gives the left-hand side.

    Alternatively, first choose $M$ in $\binom{n}{m}$ ways. Each of the remaining $n-m$ elements may independently be included in $K$ or not, giving $2^{n-m}$ choices for $K$. Therefore
    \[
        \sum_{k=m}^{n}\binom{k}{m}\binom{n}{k}
        =\binom{n}{m}2^{n-m}.
    \]

    \item Consider all subsets $A\subseteq[n]$ with $|A|\leq m$, assigning weight $(-1)^{|A|}$ to $A$. Their total weight is
    \[
        \sum_{k=0}^{m}(-1)^k\binom{n}{k}.
    \]
    Pair every $A\subseteq[n-1]$ with $|A|\leq m-1$ with $A\cup\{n\}$. The two members of each pair have opposite weights and cancel. This is a sign-reversing involution. The only unpaired subsets are the $m$-element subsets of $[n-1]$, of which there are $\binom{n-1}{m}$, each with weight $(-1)^m$. Hence
    \[
        \sum_{k=0}^{m}(-1)^k\binom{n}{k}
        =(-1)^m\binom{n-1}{m}.
    \]
\end{problemsubparts}
\end{solution}

\begin{solution}[2]
\begin{problemsubparts}
    \item By symmetry of binomial coefficients,
    \[
        \binom{n-k}{m-k}=\binom{n-k}{n-m}.
    \]
    Setting $i=n-k$ and applying Problem~2 gives
    \[
    \begin{aligned}
        \sum_{k=0}^{m}\binom{n-k}{m-k}
        &=\sum_{i=n-m}^{n}\binom{i}{n-m}\\
        &=\binom{n+1}{n-m+1}\\
        &=\binom{n+1}{m}.
    \end{aligned}
    \]

    \item By the binomial theorem,
    \[
        (1+x)^u=\sum_{k=0}^{u}\binom{u}{k}x^k,
        \qquad
        (1+x)^v=\sum_{j=0}^{v}\binom{v}{j}x^j.
    \]
    Therefore the coefficient of $x^m$ in their product is
    \[
        \sum_{k=0}^{m}\binom{u}{k}\binom{v}{m-k}.
    \]
    Since
    \[
        (1+x)^u(1+x)^v=(1+x)^{u+v},
    \]
    the same coefficient is $\binom{u+v}{m}$, proving the identity.

    \item We have
    \[
    \begin{aligned}
        \binom{k}{m}\binom{n}{k}
        &=\frac{k!}{m!(k-m)!}\frac{n!}{k!(n-k)!}\\
        &=\binom{n}{m}\binom{n-m}{k-m}.
    \end{aligned}
    \]
    Hence, with $j=k-m$,
    \[
    \begin{aligned}
        \sum_{k=m}^{n}\binom{k}{m}\binom{n}{k}
        &=\binom{n}{m}\sum_{k=m}^{n}\binom{n-m}{k-m}\\
        &=\binom{n}{m}\sum_{j=0}^{n-m}\binom{n-m}{j}\\
        &=\binom{n}{m}2^{n-m},
    \end{aligned}
    \]
    where the last equality follows from the binomial theorem.

    \item Using Pascal's identity,
    \[
        \binom{n}{k}
        =\binom{n-1}{k}+\binom{n-1}{k-1},
    \]
    we obtain
    \[
    \begin{aligned}
        \sum_{k=0}^{m}(-1)^k\binom{n}{k}
        &=\sum_{k=0}^{m}(-1)^k\binom{n-1}{k}
          +\sum_{k=1}^{m}(-1)^k\binom{n-1}{k-1}\\
        &=\sum_{k=0}^{m}(-1)^k\binom{n-1}{k}
          -\sum_{k=0}^{m-1}(-1)^k\binom{n-1}{k}\\
        &=(-1)^m\binom{n-1}{m}.
    \end{aligned}
    \]
\end{problemsubparts}
\end{solution}

\begin{remark}
Part (b) is commonly known as \emph{Vandermonde's identity}, or \emph{Vandermonde's convolution}. More generally, for $n_1,\ldots,n_r,m\in\Nzero$,
\[
    \sum_{k_1+\cdots+k_r=m}
    \binom{n_1}{k_1}\cdots\binom{n_r}{k_r}
    =\binom{n_1+\cdots+n_r}{m},
\]
with the same zero convention for out-of-range binomial coefficients. Combinatorially, this counts an $m$-element subset chosen from $r$ disjoint blocks of sizes $n_1,\ldots,n_r$.
\end{remark}

\end{document}
